% ============================================================
%  01-comandos.tex
%  Comandos y entornos propios de Terabyte para el protocolo
%  documental del taller (instrucciones_llm.md): citas con fuente
%  y pagina, etiquetado de contenido, y estructuras recurrentes.
% ============================================================

% ------------------------------------------------------------
% 1) CITAS: CP (Bases Tecnicas del Caso), BA (Bases Administrativas),
%    BTT (Bases Tecnicas Transversales)
%    Uso:  \citaCP{12}          -> (CP, p. 12)
%          \citaCP{12-14}       -> (CP, pp. 12-14)
%          \cita{BTT}{45}       -> (BTT, p. 45)
% ------------------------------------------------------------
\newcommand{\cita}[2]{\textcolor{terabyteGris}{(#1, p.\,#2)}}
\newcommand{\citaCP}[1]{\cita{CP}{#1}}
\newcommand{\citaBA}[1]{\cita{BA}{#1}}
\newcommand{\citaBTT}[1]{\cita{BTT}{#1}}

% ------------------------------------------------------------
% 2) ETIQUETAS DE CONTENIDO: Hecho / Propuesta de Terabyte /
%    Supuesto / Duda. Se usan en linea, dentro del parrafo,
%    para marcar la naturaleza de cada afirmacion.
%    Uso: El sistema opera 24/7 \Hecho. Proponemos... \Propuesta.
% ------------------------------------------------------------
\newcommand{\EtiquetaBase}[2]{%
  \textcolor{white}{\colorbox{#1}{\fuentekicker\scriptsize\bfseries\,#2\,}}%
}
\newcommand{\Hecho}{\EtiquetaBase{terabyteTeal}{HECHO}}
\newcommand{\Propuesta}{\EtiquetaBase{terabyteAcentoOscuro}{PROPUESTA TERABYTE}}
\newcommand{\Supuesto}{\EtiquetaBase{terabyteGris}{SUPUESTO}}
\newcommand{\Duda}{\EtiquetaBase{terabyteFondo}{DUDA}}

% Version en bloque (parrafo completo / bloque de analisis), con
% barra lateral de color y encabezado. Usar cuando la marca aplica
% a un bloque entero y no a una sola frase.
\NewDocumentEnvironment{bloque}{m m}{%
  % #1 = tipo (Hecho, Propuesta, Supuesto, Duda) #2 = color asociado
  \begin{tcolorbox}[
    enhanced, breakable,
    colback=white, colframe=#2,
    boxrule=0pt, leftrule=3pt,
    arc=0pt, outer arc=0pt,
    left=10pt, right=8pt, top=6pt, bottom=6pt,
    before skip=8pt, after skip=8pt,
    fonttitle=\fuentekicker\scriptsize\bfseries,
    coltitle=#2,
    colbacktitle=white,      % sin esto el titulo queda del color de su fondo
    titlerule=0pt,
    toptitle=2pt, bottomtitle=0pt,
    title={#1}
  ]
}{%
  \end{tcolorbox}
}

\NewDocumentEnvironment{bloqueHecho}{}{\begin{bloque}{HECHO}{terabyteTeal}}{\end{bloque}}
\NewDocumentEnvironment{bloquePropuesta}{}{\begin{bloque}{PROPUESTA DE TERABYTE}{terabyteAcentoOscuro}}{\end{bloque}}
\NewDocumentEnvironment{bloqueSupuesto}{}{\begin{bloque}{SUPUESTO}{terabyteGris}}{\end{bloque}}
\NewDocumentEnvironment{bloqueDuda}{}{\begin{bloque}{DUDA}{terabyteFondo}}{\end{bloque}}
% ------------------------------------------------------------
% 3) ENUNCIADO DE RIESGO
%    Estructura obligatoria: "Dado que... entonces puede que...
%    teniendo como consecuencia..."
%    Uso: \riesgo{<condicion/silencio>}{<evento incierto>}{<consecuencia>}
% ------------------------------------------------------------
\NewDocumentEnvironment{riesgo}{o}{%
  \IfValueTF{#1}{%
    \begin{tcolorbox}[
      enhanced, breakable,
      colback=terabyteCrema, colframe=terabyteAcento,
      boxrule=0.6pt, arc=2pt,
      left=10pt, right=10pt, top=7pt, bottom=7pt,
      before skip=8pt, after skip=8pt,
      fonttitle=\fuentekicker\scriptsize\bfseries,
      coltitle=terabyteFondo,
      colbacktitle=terabyteAcento!25,
      title={RIESGO #1}
    ]
  }{%
    \begin{tcolorbox}[
      enhanced, breakable,
      colback=terabyteCrema, colframe=terabyteAcento,
      boxrule=0.6pt, arc=2pt,
      left=10pt, right=10pt, top=7pt, bottom=7pt,
      before skip=8pt, after skip=8pt
    ]
  }
}{%
  \end{tcolorbox}
}

\newcommand{\enuncioriesgo}[3]{%
  \textbf{Dado que} #1, \textbf{entonces puede que} #2, \textbf{teniendo como consecuencia} #3.%
}

% ------------------------------------------------------------
% 4) CONTROL DE VERSIONES
%    Tabla estandar para la seccion formal "Control de versiones".
%    Uso:
%      \begin{controlversiones}
%        1.0 & 24-08-2026 & F. Solis / F. Nombre & Version inicial \\
%      \end{controlversiones}
% ------------------------------------------------------------
\newenvironment{controlversiones}{%
  \begin{center}
  \textotabla
  \begin{tabular}{@{}p{1.6cm}p{2.2cm}p{4.0cm}p{\dimexpr\textwidth-1.6cm-2.2cm-4.0cm-8\tabcolsep\relax}@{}}
    \toprule
    \rowcolor{terabyteFondo}
    \textcolor{white}{\bfseries Versión} &
    \textcolor{white}{\bfseries Fecha} &
    \textcolor{white}{\bfseries Responsable(s)} &
    \textcolor{white}{\bfseries Cambios principales} \\
    \midrule
}{%
    \bottomrule
  \end{tabular}
  \end{center}
}

% ------------------------------------------------------------
% 5) TARJETAS NUMERADAS (equivalente a los bloques "01/02/03" de la
%    linea grafica, utiles para resumenes de tres o cuatro puntos)
%    Uso:
%      \begin{tarjetas}{3}
%        \tarjeta{Capacidad tecnica}{Descripcion breve.}
%        \tarjeta{Experiencia comparable}{Descripcion breve.}
%        \tarjeta{Equipo asignado}{Descripcion breve.}
%      \end{tarjetas}
% ------------------------------------------------------------
\newcounter{tarjetaidx}
\NewDocumentEnvironment{tarjetas}{m}{%
  \setcounter{tarjetaidx}{0}%
  \begin{center}
  \renewcommand{\arraystretch}{1}%
}{%
  \end{center}
}
\newcommand{\tarjeta}[2]{%
  \stepcounter{tarjetaidx}%
  \begin{minipage}[t]{0.31\textwidth}
    {\color{terabyteAcento}\rule{\linewidth}{2pt}}\\[4pt]
    {\fuentekicker\scriptsize\color{terabyteAcentoOscuro}\arabic{tarjetaidx}\,/\,03}\\[2pt]
    {\fuentetitulo\normalsize\color{terabyteFondo}#1}\\[4pt]
    {\small\color{terabyteTextoCuerpo}#2}
  \end{minipage}\hfill
}

% ------------------------------------------------------------
% 6) TITULOS DE SECCION EN LINEA CON EL ESTILO CORPORATIVO
% ------------------------------------------------------------
\titleformat{\section}
  {\fuentetitulo\Large\color{terabyteFondo}}
  {\thesection.}{0.5em}{}
  [\vspace{2pt}{\color{terabyteAcento}\titlerule[1.4pt]}]
\titlespacing*{\section}{0pt}{22pt}{12pt}

\titleformat{\subsection}
  {\fuentetitulo\large\color{terabyteFondo}}
  {\thesubsection}{0.6em}{}
\titlespacing*{\subsection}{0pt}{16pt}{8pt}

\titleformat{\subsubsection}
  {\fuentetitulo\normalsize\color{terabyteAcentoOscuro}}
  {\thesubsubsection}{0.6em}{}
\titlespacing*{\subsubsection}{0pt}{12pt}{6pt}

% "Kicker" de apertura de seccion (etiqueta pequena tipo
% "02 · PROBLEMA Y NECESIDAD" vista en la portada/slides)
\newcommand{\kicker}[1]{{\fuentekicker\small\color{terabyteAcentoOscuro}\MakeUppercase{#1}}\par\vspace{2pt}}

% ------------------------------------------------------------
% 7) ENCABEZADOS Y PIES DE PAGINA
% ------------------------------------------------------------
\newcommand{\nombreempresa}{Terabyte}
\newcommand{\tituloactual}{}
\newcommand{\establecertitulo}[1]{\renewcommand{\tituloactual}{#1}}

\pagestyle{fancy}
\fancyhf{}
\renewcommand{\headrulewidth}{0.4pt}
\renewcommand{\footrulewidth}{0pt}
\renewcommand{\headrule}{{\color{terabyteLineaClara}\hrule height\headrulewidth}}
\fancyhead[L]{\fuentekicker\scriptsize\color{terabyteAcentoOscuro}\MakeUppercase{\leftmark}}
\fancyhead[R]{\fuentekicker\scriptsize\color{terabyteGris}\tituloactual}
\fancyfoot[C]{%
  \begin{tikzpicture}[remember picture, overlay]
    \fill[terabyteAcento] ([yshift=0.55cm]current page.south west) rectangle ([yshift=0.62cm]current page.south east);
  \end{tikzpicture}%
  \fuentekicker\scriptsize\color{terabyteGris}\thepage
}

\fancypagestyle{plain}{%
  \fancyhf{}
  \renewcommand{\headrulewidth}{0pt}
  \fancyfoot[C]{%
    \begin{tikzpicture}[remember picture, overlay]
      \fill[terabyteAcento] ([yshift=0.55cm]current page.south west) rectangle ([yshift=0.62cm]current page.south east);
    \end{tikzpicture}%
    \fuentekicker\scriptsize\color{terabyteGris}\thepage
  }
}

% \leftmark en formato prolijo: "01 · TITULO" en el cuerpo numerado,
% "ANEXO A · TITULO" en los anexos (letras).
\newif\ifenanexo
\renewcommand{\sectionmark}[1]{%
  \ifenanexo
    \markboth{ANEXO~\thesection\ \textperiodcentered\ \MakeUppercase{#1}}{}%
  \else
    \markboth{\ifnum\value{section}<10 0\fi\thesection\ \textperiodcentered\ \MakeUppercase{#1}}{}%
  \fi
}
